% Solutions to the Chapter 3 exercises that are not code, from lectures/L03/appendix/c_exercises.md.
% Exercises 3.6 and 3.7 are Code and are not answered; see chapter.tex for why.

\section{Chapter 3: The Kernel: Architecture, Source Tree, and Configuration}
\label{sol:sec:c3}
Answers to the exercises in \secref{c3:sec:exercises}. Exercises~\ref{c3:ex:6} and~\ref{c3:ex:7}
are Code: the first is checked by \code{make test L=L03}, and the second is answered with a command
and a number on your own target.

% ------------------------------------------------------------------------------------------
\solution[fillaccent2]{c3:ex:1}{What kind of program the kernel is}{Recall}

\xpart{a} Any three of these (\secref{c3:app:a1}):
\begin{itemize}
  \item \textbf{It starts at \code{main} and is over when \code{main} returns.} The kernel has no
    \code{main} and never returns; after boot it runs only when something enters it: a system call,
    an interrupt, a fault.
  \item \textbf{It has an address space that nothing else can write.} The kernel is one address
    space shared by all of its parts, so a bug in any driver can corrupt any subsystem.
  \item \textbf{A crash kills it and nothing else.} A process that faults is cleaned up by the
    kernel; when the kernel faults, there is nothing above it to clean up.
  \item \textbf{It is linked against a C library.} The kernel is built \code{-nostdinc} and
    \code{-ffreestanding}, against its own headers.
  \item \textbf{Its stack grows on demand.} The kernel's is fixed: 16 KB per task on arm64.
\end{itemize}

\xpart{b} It claims that the whole kernel, core, drivers, filesystems and network stack alike, runs
in one privileged address space with no protection between its parts: any of them can read and
write the others' memory, and a fault in one is a fault in the kernel. It does not claim that the
kernel is one indivisible piece, since modules are loaded and unloaded at run time, into that same
address space. And it says nothing about how the source is organised, which is intensely modular,
with a directory, maintainers and internal interfaces for each subsystem. Monolithic is a claim
about the address space, not about modularity (\secref{c3:app:a1}); the alternative it is contrasted
with is a microkernel, where drivers and filesystems run as separate processes that the hardware
keeps apart.

\xpart{c} An oops. The fault handler finds no fixup for the address, prints \code{Unable to handle
kernel NULL pointer dereference at virtual address 0000000000000000} with a register dump and a
backtrace, taints the kernel, and kills the task that was running. If there is no task to blame,
because the fault was in an interrupt handler, it panics instead, and it panics anyway if
\code{panic_on_oops} is set. The difference is structural. A process's fault is handled by a
separate program that is still intact, the kernel, which owns everything the process held and can
take all of it back. When the kernel faults, the code that would clean up is the code that faulted,
in the address space that may now be corrupt, and nobody is above it. Killing the task does not
release the locks it held or finish the update it was halfway through, so even an oops the machine
survives leaves it in an unknown state. Better testing means fewer faults; it cannot give the kernel
a supervisor.

\xpart{d} \code{printf} belongs to the C library, and the C library is built on top of the kernel:
\code{printf} formats into a buffer and hands it to \code{write()}, which is a system call, a request
to the kernel. The kernel cannot make that request of itself, and it is built without any C library
at all. What it has instead is \code{printk}, which formats into the kernel's own log buffer, the
one \code{dmesg} reads, and out to the console, carries a log level, and can be called from any
context, an interrupt handler included (\secref{c4:app:a7}).

% ------------------------------------------------------------------------------------------
\solution[fillaccent2]{c3:ex:2}{The two interfaces}{Recall}

\xpart{a} The system call interface is stable forever: a change that makes a working userspace
program stop working is reverted, however correct it was. The internal API has no guarantee at all:
function signatures, structures and whole subsystems may change in any release
(\secref{c3:app:a2}).

\xpart{b} \textbf{From the stable system call interface:} the userspace half of your product, the
programs that talk to your driver through \code{open}, \code{read} and \code{ioctl} and the rest of
the userland around them, keeps working across kernel upgrades and never needs rebuilding for one.
The same rule applies to what your driver adds: an \code{ioctl} or a sysfs file you ship is expected
to stay as it is for good. \textbf{From the unstable internal API:} every new kernel may break your
driver's build, or worse, keep it building and change what a function it calls means. You rebuild
and retest against every kernel you support, port the driver whenever an interface it uses moves,
and hear about the change from the compiler, because nobody making it knows your driver exists.

\xpart{c} A stable internal ABI would freeze the kernel's internal design, mistakes included, for
good. The kernel improves, and fixes its security bugs, by changing an internal interface and every
caller of it in the same patch; with binary drivers depending on the old one, nothing they touch
could change again. And a binary interface is not even one thing to freeze: the layout of the same
structure differs with the compiler version, the configuration (SMP, the preemption model, debugging
options) and the architecture. The file is \file{Documentation/process/stable-api-nonsense.rst},
which makes both arguments and ends with the answer the project prefers: put the driver in the tree.

\xpart{d} Whoever changes the API. The tree must build at every commit, so a patch that changes an
internal interface converts every caller in the tree in the same series, usually mechanically and
often with Coccinelle, or the maintainer will not take it. Your driver is one of those callers, and
you are typically asked to review the change rather than to make it. That is the maintenance cost
an out-of-tree driver pays and an in-tree one does not.

% ------------------------------------------------------------------------------------------
\solution[fillaccent3]{c3:ex:3}{Reading Kconfig}{Hand calculation}

Both entries sit inside the \code{if GPIOLIB} block of \file{drivers/gpio/Kconfig}, so all of this
assumes \code{CONFIG_GPIOLIB=y}, as \file{kernel/qa.config} sets it.

\xpart{a} No \code{GPIO_SYSFS} line at all, and \code{CONFIG_GPIO_CDEV=y}. With \code{EXPERT} off,
\code{GPIO_SYSFS}'s prompt is invisible, and a value from a fragment is only ever applied to a
symbol that has a visible prompt, so the \code{y} is ignored. The symbol has no default either, so
its value is \code{n}, and a symbol that is neither visible nor given a default is not written to
the \code{.config}: not even as \code{\# CONFIG_GPIO_SYSFS is not set}. \code{GPIO_CDEV} keeps its
\code{default y}. The course's \file{ci/kernel.sh} notices the line that vanished and prints
\code{warning: CONFIG_GPIO_SYSFS=y did not survive olddefconfig} (\secref{c3:app:b2}).

\xpart{b} \code{CONFIG_EXPERT=y}, \code{CONFIG_GPIO_SYSFS=y} and \code{CONFIG_GPIO_CDEV=y}. The
prompt is now visible, its \code{depends on SYSFS} is met, so the fragment's \code{y} is taken; and
the \code{select} forces \code{GPIO_CDEV} to \code{y}, which it already was by default, but now it
can no longer be turned off while \code{GPIO_SYSFS} is on. \code{EXPERT} uncovers many other
prompts too, and \code{olddefconfig} answers each of them with its default.

\xpart{c} Yes, \code{CONFIG_GPIO_CDEV=y}. A symbol whose prompt is invisible is never asked about,
but it is still evaluated: its value comes from its default, here \code{y}, and because it has one
it is written out. Visibility decides whether you are asked; the value of a symbol nobody was asked
about comes from its \code{default}, or from a \code{select} by a symbol that is on. That is why
this course's \code{.config} has \code{CONFIG_GPIO_CDEV=y} with \code{EXPERT} off.

\xpart{d} \code{depends on B} is a precondition: until \code{B} holds, the symbol is hidden and
held at \code{n}. \code{select B} is an order: while the selecting symbol is on, \code{B} is forced
on, without looking at \code{B}'s own dependencies. \textbf{\code{select}} is the one that can break
a \code{depends on}: if \code{A} selects \code{B} and \code{B} depends on \code{C}, then turning on
\code{A} with \code{C} off forces \code{B} on anyway, and the \code{.config} holds a \code{B} whose
\code{depends on C} is false. Kconfig prints \code{WARNING: unmet direct dependencies detected for
B}, and the build may fail inside \code{B}'s code, which was written on the assumption that \code{C}
is there. That is also how an option nobody selected, with an unsatisfied \code{depends on}, ends
up in a \code{.config}.

\xpart{e} \code{CONFIG_GPIO_CDEV}, and it works, although the line is not what makes it work: with
\code{EXPERT} off the symbol has no prompt, so the fragment's value is ignored exactly as in
\textbf{a)}, and it is \code{y} anyway because that is its default. The line survives
\code{olddefconfig} only because it agrees with the default; the check in \file{ci/kernel.sh}
passes, and that is what this course's \file{kernel/qa.config} does. The difference shows the
other way round: \code{\# CONFIG_GPIO_CDEV is not set} would not take effect either, so with
\code{EXPERT} off you cannot turn it off at all. (For a \file{/dev/gpiochipN} to appear there must
also be a GPIO controller to represent; this target has one, the PL061.)

% ------------------------------------------------------------------------------------------
\solution[fillaccent3]{c3:ex:4}{Built in or module}{Hand calculation}

\xpart{a} Must be \code{=y}: with no initramfs the kernel mounts the root filesystem itself, and a
module lives on the root filesystem it would be needed to reach.

\xpart{b} May be \code{=m}: the initramfs is a root filesystem in RAM that the bootloader loads
with the kernel, and its init loads the driver from there before mounting the real root.

\xpart{c} Cannot be \code{=m}: a \code{bool} has only the values \code{y} and \code{n}, so it is
\code{=y} if anything on the system uses GPIO and \code{=n} otherwise.

\xpart{d} May be \code{=m}: it is needed only while the device is plugged in, it is loaded then
through its module alias, and until then it costs nothing.

\xpart{e} Must be \code{=y}, for the same reason as \textbf{a)}: the kernel has to read the root
filesystem before it can load any module, so the code that reads it cannot be one.

\xpart{f} Cannot be \code{=m}: \code{PREEMPT_RT} is a \code{bool} in the preemption-model choice
and changes how the whole kernel is compiled, turning spinlocks into sleeping locks and interrupt
handlers into threads, which is not something that can be loaded later. It is \code{=y} only in
\file{build/kernel-rt}, and only because \file{kernel/rt.config} turns \code{EXPERT} on as well.

% ------------------------------------------------------------------------------------------
\solution[fillaccent4]{c3:ex:5}{Navigating the tree}{Design}

\xpart{a} \file{kernel/sched/}. \file{core.c} holds \code{__schedule()} and
\code{pick_next_task()}, which asks each scheduling class in turn; \file{fair.c} is the class almost
every task is in. If the name is not known: \code{grep -rn "pick_next_task" kernel/sched/}.

\xpart{b} \file{drivers/hwmon/}: a temperature sensor is a hardware monitor, and most I2C ones are
there (\file{lm75.c}, \file{tmp102.c}). One used for measurement rather than monitoring may be an
IIO driver in \file{drivers/iio/temperature/}. Search by the chip's name, which is what the
datasheet gives you: \code{grep -rli "tmp102" drivers/}.

\xpart{c} \file{include/linux/}: it is in \file{fs.h}, at line 2062 of the pinned tree. To find the
definition among thousands of uses, anchor the search on the opening line: \code{grep -rn "^struct
file_operations \{" include/}.

\xpart{d} \file{arch/arm64/kernel/}: \file{entry.S} holds the exception vector table
(\code{vectors}) and the \code{kernel_ventry} macro, and \file{entry-common.c} the C handlers it
calls, such as \code{el0t_64_sync_handler} for a system call or a fault from userspace. Search:
\code{grep -rn "SYM_CODE_START(vectors)" arch/arm64/}.

\xpart{e} \file{fs/}: \file{open.c} defines it as \code{SYSCALL_DEFINE4(openat, ...)}, at line 1441,
and it calls \code{do_sys_open()}. Search for the macro, because the name \code{sys_openat} is
generated by it and appears in no definition: \code{grep -rn "SYSCALL_DEFINE.(openat," fs/}.

\xpart{f} \file{kernel/}, because it is about process 1 exiting, which is process management.
\code{grep -rn "Attempted to kill init" kernel/} finds \file{kernel/exit.c}, line 908, in a call to
\code{panic()}. The format string goes on \code{exitcode=0x\%08x}, so search for the fixed words and
never for the number the message printed.

% ------------------------------------------------------------------------------------------
\solution[fillaccent4]{c3:ex:8}{A configuration for a real product}{Design}

\xpart{a} Seven components, seven decisions:
\begin{itemize}
  \item \textbf{The SoC: \code{=y}.} Its platform symbol, interrupt controller, clocks, pin control
    and timer are needed before anything else can run, and most of them are \code{bool} anyway.
    Every other SoC family that \code{defconfig} enables is \code{=n}.
  \item \textbf{The 256 MB of RAM: nothing to choose.} Memory management is core kernel, always
    built in. What the size decides is how much you can afford to leave switched on, not whether
    anything is a module.
  \item \textbf{Ethernet: \code{=y}.} It is used on every boot, and the field updates of the sensor
    driver presumably arrive over it; as a module it would only add a step that can fail. \code{=m}
    is defensible if image size is tight, since nothing needs the network before userspace; that
    argument has to show the space matters.
  \item \textbf{The eMMC: \code{=y}}, meaning the MMC core, the SoC's host controller driver and the
    MMC block driver: the root filesystem is on it and there is no initramfs to load a module from
    (Exercise~\ref{c3:ex:4}\textbf{a)}).
  \item \textbf{ext4: \code{=y}}, for the same reason (Exercise~\ref{c3:ex:4}\textbf{e)}).
  \item \textbf{The I2C sensor: \code{=m}.} Replacing a driver without a reboot is only possible
    for a module: install the new \code{.ko}, \code{rmmod} the old one, load the new. The I2C core
    and the SoC's I2C controller driver are \code{=y}: they are not the part that is updated, and
    the sensor module needs them.
  \item \textbf{The watchdog: \code{=y}.} If U-Boot has already started it, the kernel has to keep
    it fed from early boot until the userspace daemon takes over, which a module loaded later
    cannot do; and the driver is small.
\end{itemize}

\xpart{b} The ones on the path to the root filesystem: the eMMC driver, including anything its
host controller needs from the SoC such as clocks and pin control, and ext4. Either one as
\code{=m} gives the same symptom: the kernel boots, runs its initialisation, and cannot mount the
root filesystem. With no eMMC driver there is no block device, so it prints \code{VFS: Cannot open
root device}, lists the partitions it can see, which are none, and stops with \code{Kernel panic -
not syncing: VFS: Unable to mount root fs on unknown-block(0,0)}. With the device present but no
ext4, it prints \code{No filesystem could mount root, tried:} and the types it does have, and
panics with the same message naming the device it found. One more decision stops the boot if
U-Boot has started the watchdog: with the driver \code{=n}, or \code{=m} and loaded too late,
nothing feeds it, and the board resets partway through every boot.

\xpart{c} Keep one kernel for both variants. Build both sensor drivers as modules, which they are
already, and let each variant's device tree say which sensor is fitted, so the right driver is
loaded by its \code{compatible} string; the variants then differ in their device tree, not in their
kernel. \code{savedefconfig} is how the product's configuration is stored: after any change, it
writes a \code{defconfig} holding only the choices that differ from the defaults, which is small
enough to review in a commit, readable in a diff, and brought up to date by \code{olddefconfig} when
the kernel moves on. If a variant ever does need a different kernel option, keep it as a fragment
merged over that base, the way \file{kernel/rt.config} sits over this course's configuration, rather
than as a second full \code{.config}.

\xpart{d} The configuration: how many modules it builds (\file{modules.order} has one line per
module) and which hardware it builds them for. Expect to find a vendor or architecture
\code{defconfig} building for every board anyone upstreamed: dozens of SoC families with their
clock, pin control, DMA and PHY drivers, graphics, sound and wireless, thousands of modules that can
never load on the product. Turning those off is what took this course's build from about an hour
and roughly two thousand modules to eighteen minutes and 458 (\secref{c3:app:b7}). The same file
turns off DWARF debugging information, which by its own comment roughly doubles the build time.

% ------------------------------------------------------------------------------------------
\solution[fillaccent]{c3:ex:9}{What one config option costs}{Cross-check}

\xpart{a} $76 + 34{,}048 + 10 + 4 = 34{,}138$ bytes, from \code{.text}, \code{.rodata},
\code{.rodata.str1.8} and \code{.initcall6.init} (\secref{c3:app:b6}). The listing's \code{Total},
34,532, also counts the metadata and the 148 bytes of \code{__init} and \code{__exit} code.

\xpart{b} They should not be identical, and each difference is accounted for exactly.
\textbf{\code{configs.o}'s 34,048} is more than the compressed file: \file{kernel/configs.c} wraps
the blob in two 8-byte markers, \code{IKCFG_ST} and \code{IKCFG_ED}, which
\file{scripts/extract-ikconfig} searches for, and the same section holds the 96-byte
\code{proc_ops} table behind \file{/proc/config.gz}; $33{,}936 + 16 + 96 = 34{,}048$. \textbf{The
host's 33,944} includes the file name: \code{gzip} records \code{.config} and its terminating zero
byte, 8 bytes, in the header, while the kernel's build compresses with \code{gzip -n}, which leaves
the name out; $33{,}936 + 8 = 33{,}944$. \textbf{The target's 33,936} is the blob and nothing else,
because \code{ikconfig_init} sets the file's size to the distance between the two symbols that
bracket it; it is the size of \file{build/kernel/kernel/config_data.gz}.

\xpart{c} The course's baseline, published in \secref{c3:app:a8}: \code{Image} 24,943,104 bytes,
\code{vmlinux} 31,468,440 bytes, and among the sections \code{.text} at 13,021,184 and
\code{.rodata} at 4,369,606.

\xpart{e} The course measured two of the three and published them, drawn in
Figure~\ref{fig:sol:c3:ikconfig}: \code{vmlinux}'s section total fell by 41,208 bytes, and the
\code{Image} by 65,536. It did not publish the change in the \code{vmlinux} file; \textbf{f)} says
what decides it.

\lecturefigure[0.75]{exam/images/ikconfig_sizes.png}
  {The one-line change of Exercise~\ref{c3:ex:9}, measured three ways on the course's target: 34,138
   bytes of content in \code{configs.o}, 41,208 bytes of \code{vmlinux} sections, 65,536 bytes of
   \code{Image}. The first gap is page alignment and the second 64 KiB segment alignment.}
  {fig:sol:c3:ikconfig}

\xpart{f} Each of the four:
\begin{itemize}
  \item \textbf{Closest and furthest.} The section total is closest, at $41{,}208 / 34{,}138 =
    1.21$ times the prediction. The \code{Image} is $65{,}536 / 34{,}138 = 1.92$ times it, and the
    \code{vmlinux} file is further still. \code{readelf -l} on \file{build/kernel/vmlinux} shows why:
    the file's loadable part is the \code{Image} byte for byte, starting 64 KiB into the file, and
    the symbol and string tables follow it. So the file shrinks by the \code{Image}'s 65,536 plus 24
    bytes of \code{.symtab} for each symbol \code{configs.o} defined and the bytes of their names in
    \code{.strtab}: a little more than the \code{Image}, and the furthest of the three, by a factor
    of just over 1.9.
  \item \textbf{The per-section deltas.} \code{.rodata} fell by 36,864 bytes for the 34,058 that
    \code{configs.o} put there ($34{,}048 + 10$), and \code{.text} by 4,096 for its 76. Each delta
    is a whole number of 4,096-byte pages, the content rounded up to the page: $34{,}058 / 4{,}096
    = 8.3$, so nine pages, 36,864; 76 bytes, one page. Both sections hold page-aligned material
    behind \code{configs.o}'s bytes (in \code{.text}, the hypervisor code, which starts and ends on
    a page), so removing anything moves it down in whole pages. Which way the rounding goes depends
    on how much padding was already in front of the boundary: 76 bytes moved \code{.text} by a page
    only because the code before the boundary spilled less than 76 bytes into its last page; on
    another build the same 76 bytes could move it by nothing. The other three deltas are small:
    \code{.init.text} fell by 112 for \code{configs.o}'s 108, rounded to 8-byte alignment;
    \code{.exit.text} by exactly its 40; and \code{.rela.dyn} by 96, four 24-byte relocation records
    the linker kept for addresses stored in \code{configs.o}'s data, which the kernel patches when
    it relocates itself at boot. $36{,}864 + 4{,}096 + 112 + 96 + 40 = 41{,}208$.
  \item \textbf{The \code{Image} delta} is $65{,}536 = 64$~KiB exactly, one unit of
    \code{SEGMENT_ALIGN}, which is \code{SZ_64K} in \file{arch/arm64/include/asm/memory.h}. The
    linker script aligns the end of the text, the start of the init segment, the boundary between
    init text and init data, and the end of the init segment to it. So the \code{Image} is a row of
    segments, each starting on a 64 KiB boundary with padding in front of it, which lets each be
    mapped with its own permissions using the largest mappings the hardware offers. The
    \code{Image} can therefore only change size in 64 KiB steps, and a step happens when some
    segment's contents cross a boundary.
  \item \textbf{A second 34 KB option} would not necessarily shrink the \code{Image} by 64 KiB again.
    About 41 KB of sections is less than one step, so it moves the \code{Image} by 65,536 only if it
    takes a segment's contents back across a boundary, which depends on how much slack the first
    removal left in front of each one: it could be 65,536 again, or nothing. Over many options the
    \code{Image} shrinks by roughly the sum of what they contain; any single measurement is
    quantised.
\end{itemize}

\xpart{g} A report that handles both worlds passes again, with seventeen checks: \file{run.sh} says
the configuration is available as \code{inferred}, and the report reads \code{config-source:
inferred} and \code{hz: unknown}.

\xpart{h} ``About 34 KB: the compressed copy of the configuration and the code that serves it as
\file{/proc/config.gz}, which the kernel keeps in memory for as long as it runs.'' That is the
content, 34,138 bytes, and it is the one to quote because it is the option's own cost and the same
on any build. The other two measure padding as well, which belongs to wherever the boundaries
happened to fall on this build: the same option could cost 64 KiB of \code{Image} here and nothing
on the next build. If the colleague is fitting this particular \code{Image} into a flash partition,
64 KiB is what they pay on this build; then say that it is a number about this build, not about the
option.
